The registered load probe completes 85/85 valid schedules, of which 68 meet their deadlines, in batches of 1, 4, 16 and 64 with at most four concurrent checker processes. Observed throughput rises from 21.06 to 37.07, 51.09 and 56.95 tasks per second; the 95th-percentile enqueue-to-durable latency rises from 37.63 to 89.63, 288.45 and 1,018.46\,ms. No batch exhausts its deadline, and no mid-batch manual intervention occurs. These are finite local measurements, including queueing, input preparation, evidence writes and checker startup/shutdown. They do not establish sustained throughput under arbitrary load.

The complete confirmatory campaign uses one explicit launch for 336 stream/treatment executions and 72 native reference calls, with zero manually triggered intermediate phases. It takes 844.40 seconds of elapsed time and 751.10 seconds of measured Node CPU; the separate Python CPU total is unavailable. Common synthetic WORLD construction takes 420.61\,ms and includes 3,464,226 known baseline-generation units, reported separately from conditional per-arm mission work. Setup, authoring, review and the additional historical replay remain research effort, not free autonomous output.

The 336 completed trial directories occupy 209,075,637 bytes in total, with final checkpoints ranging from 89,913 to 442,729 bytes. In the primary flywheel arm, checkpoint serialization alone costs 5,238,180 of 11,714,786 charged units (44.71\%). Retention cannot be evaluated as though its history were free. These full research snapshots are not an optimal archival design; an incrementally stored journal could change this cost boundary and would require its own measured comparison.
